% Fuente: Auxiliar 3, sección «Forma general y estándar».
{\color{MidnightBlue}
\textbf{Pasar a forma estándar:}\\

Definimos los siguientes cambios de variables:
\[
x_3 = x_3' + 5,\quad x_3'\ge0;
\qquad
x_4 = -x_4',\quad x_4'\ge0;
\qquad
x_5 = x_5^+ - x_5^-,\quad x_5^+, x_5^-\ge0.
\]

Definimos la variable de holgura $s_1$ y se añade a la primera restricción:
\[
x_1 + x_2 +4x_4' + s_1 =10,\quad s_1\ge0;
\]
donde, además, se añade el cambio de variable de $x_4'$.\\

Se realizan los cambios de variables en la segunda restricción:
\[
x_2 -0.5x_3' + x_5^+ - x_5^- =1.5
\]

La función objetivo se pasa a minimización y se realizan los cambios de variables en la función objetivo; además, se omiten las constantes:

\[
\min\;-3x_1 -5x_2 +3x_3' +1.3x_4' +x_5^+ - x_5^-.
\]

El problema queda:
\[
\begin{aligned}
\min\quad & -3x_1 -5x_2 +3x_3' +1.3x_4' +x_5^+ - x_5^- \\
\text{s.a.}\quad
& x_1 + x_2 +4x_4' + s_1 =10,\\
& x_2 -0.5x_3' + x_5^+ - x_5^- =1.5,\\
& x_1, x_2, x_3', x_4', x_5^+, x_5^-, s_1 \ge0.
\end{aligned}
\]

Las matrices están dadas por:
\[
x = \begin{pmatrix}
x_1\\ x_2\\ x_3'\\ x_4'\\ x_5^+\\ x_5^-\\ s_1
\end{pmatrix},
\quad
c^T = \bigl(-3,\,-5,\,3,\,1.3,\,1,\,-1,\,0\bigr),
\]
\[
A =
\begin{pmatrix}
1 & 1   & 0    & 4 & 0 & 0 & 1\\
0 & 1   & -0.5 & 0 & 1 & -1 & 0
\end{pmatrix},
\quad
b = \begin{pmatrix}10\\1.5\end{pmatrix}.
\]\\

\textbf{Pasar a forma general:}\\

La segunda restricción se puede escribir como dos restricciones de desigualdad:
\[
x_2 - 0.5x_3 + x_5 = -1
\;\Longrightarrow\;
\begin{cases}
x_2 - 0.5x_3 + x_5 \le -1,\\
-(x_2 - 0.5x_3 + x_5)\le 1.
\end{cases}
\]\\

La tercera restricción se puede cambiar la dirección de la desigualdad:
\[
x_3 \ge 5
\;\Longrightarrow\;
-\,x_3 \le -5.
\]

El problema queda:
\[
\begin{aligned}
\min\quad & -3x_1 -5x_2 +3x_3 -1.3x_4 +x_5 \\
\text{s.a.}\quad
& x_1 + x_2 -4x_4 \le 10,\\
&x_2 - 0.5x_3 + x_5 \le -1,\\
&-x_2 + 0.5x_3 - x_5\le 1. \\
& -x_3 \le -5
\end{aligned}
\]
\[
x_1\ge0,\quad x_2\ge0,\quad x_4\le0,\quad x_5\ \text{libre}.
\]

Las matrices están dadas:
\[
x = \begin{pmatrix}
x_1\\ x_2\\ x_3\\ x_4\\ x_5\\ 
\end{pmatrix},
\quad
c^T = \bigl(-3,\,-5,\,3,\,-1.3,\,1),
\]
\[
A = 
\begin{pmatrix}
1 & 1   & 0    & -4 & 0\\
0 & 1   & -0.5 & 0  & 1\\
0 & -1  & 0.5  & 0  & -1\\
0 & 0   & -1   & 0  & 0
\end{pmatrix},\quad
b = \begin{pmatrix}10\\-1\\1\\-5\end{pmatrix}.
\]

\textbf{Forma estándar:}\\

El problema se encuentra en forma estándar, de modo que sólo queda representar de forma matricial:

\[
x = \begin{pmatrix}
x_1\\
x_2\\
r_1\\
s_1
\end{pmatrix},
\quad
c^T = \bigl(200,\;220,\;0,\;0\bigr),
\]
\[
A = 
\begin{pmatrix}
0.7 & 0.6 & -1 & 0\\
0.9 & 0.8 &  0 & 1\\
1   & 1   &  0 & 0
\end{pmatrix},
\quad
b = \begin{pmatrix}
0.65\\
0.85\\
1
\end{pmatrix}.
\]

\textbf{Forma general:}\\

Cada restricción se convierte en dos desigualdades “\(\le\)”:
\[
\begin{aligned}
&0.7x_1 + 0.6x_2 - r_1 \le 0.65,\\
&-(0.7x_1 + 0.6x_2 - r_1)\le -0.65,\\[6pt]
&0.9x_1 + 0.8x_2 + s_1 \le 0.85,\\
&-(0.9x_1 + 0.8x_2 + s_1)\le -0.85,\\[6pt]
&x_1 + x_2 \le 1,\\
&-(x_1 + x_2)\le -1.
\end{aligned}
\]\\

Las matrices están dadas por:
\[
x = \begin{pmatrix}
x_1\\
x_2\\
r_1\\
s_1
\end{pmatrix},\quad
c^T = \bigl(200,\;220,\;0,\;0\bigr).
\]
\[
A = 
\begin{pmatrix}
 0.7 &  0.6 & -1 &  0\\
-0.7 & -0.6 &  1 &  0\\
 0.9 &  0.8 &  0 &  1\\
-0.9 & -0.8 &  0 & -1\\
 1   &  1   &  0 &  0\\
-1   & -1   &  0 &  0
\end{pmatrix},
\quad
b = 
\begin{pmatrix}
 0.65\\
-0.65\\
 0.85\\
-0.85\\
 1\\
-1
\end{pmatrix}.
\]
}

